\begin{exo}
	On rappelle ci-dessous la composition d'un jeu de $32$ cartes.

	\begin{center}
		\def\arraystretch{1.6}
		\begin{tabular}{|c|c|c|c|c|c|c|c|} \hline
			$7 \spadesuit$           & $8\spadesuit$            & $9\spadesuit$            & $10\spadesuit$            & $J\spadesuit$            & $Q\spadesuit$            & $K\spadesuit$            & $A\spadesuit$            \\ \hline
			\color{red}{$7\diamondsuit$} & \color{red}{$8\diamondsuit$} & \color{red}{$9\diamondsuit$} & \color{red}{$10\diamondsuit$} & \color{red}{$J\diamondsuit$} & \color{red}{$Q\diamondsuit$} & \color{red}{$K\diamondsuit$} & \color{red}{$A\diamondsuit$} \\ \hline
			\color{red}{$7\heartsuit$}   & \color{red}{$8\heartsuit$}   & \color{red}{$9\heartsuit$}   & \color{red}{$10\heartsuit$}   & \color{red}{$J\heartsuit$}   & \color{red}{$Q\heartsuit$}   & \color{red}{$K\heartsuit$}   & \color{red}{$A\heartsuit$}   \\ \hline
			$7 \clubsuit$            & $8\clubsuit$             & $9\clubsuit$             & $10\clubsuit$             & $J\clubsuit$             & $Q\clubsuit$             & $K\clubsuit$             & $A\clubsuit$             \\ \hline
		\end{tabular}
	\end{center}

	Dans chaque cas, déterminer l'univers associé à l'expérience aléatoire décrite.
	\begin{enumerate}
		\item On tire une carte au hasard et on s'intéresse à sa couleur.
		\item On tire une carte au hasard et on s'intéresse à sa valeur.
	\end{enumerate}
\end{exo}
